\section{Other Renormalization Schemes}
\label{sec:ren}

Although we derive the above matching formula assuming that the quasi-PDF is renormalized in the $\overline{\ensuremath{\operatorname{MS}}}$ scheme, this is not a limitation to our result.
Since the gauge-invariant Wilson line operator $\tilde{O}_\Gamma(z)$ has been proven to be multiplicatively renormalizable in the coordinate space~\cite{Ji:2017oey,Ishikawa:2017faj}, one can convert $\tilde{Q}_{\Gamma}(z)$ from any other scheme to the $\overline{\ensuremath{\operatorname{MS}}}$ scheme before using the above factorization formula. The renormalization of the
quasi-PDF has been studied in many recent papers~\cite{Ji:2015jwa,Ishikawa:2016znu,Chen:2016fxx,Xiong:2017jtn,Constantinou:2017sej,Alexandrou:2017huk,Chen:2017mzz,Green:2017xeu,Stewart:2017tvs,Wang:2017eel}. We will discuss some
of these results and show how they can be incorporated into the factorization formula in Eq.~(\ref{eq:fac-tq}).

The $\overline{\ensuremath{\operatorname{MS}}}$ scheme is convenient for our discussion of the OPE as it guarantees Lorentz and gauge invariances, but it is not practical for lattice renormalization. Since the lattice theory has a natural UV cut-off $1/a$ with $a$ being the lattice spacing, the unrenormalized \spcorr $\tilde{Q}$ inherits the power divergence from the Wilson line self-energy according to Eq.~(\ref{eq:ren}). For an arbitrary scheme $X$, the renormalized \spcorr
\begin{align}
\tilde{Q}^X(\zeta,z^2\mu^2_R)
  = \lim_{a\to0} Z^{-1}_X(z^2\mu_R^2,a^2\mu^2_R)\, \tilde{Q}(\zeta,z^2/a^2)
\end{align}
should be free of all the UV divergences and have a well-defined continuum limit as $a\to0$. This continuum limit, in particular, is independent of the UV regulator, so
\begin{align}
&\lim_{a\to0} Z^{-1}_X(z^2\mu_R^2,a^2\mu^2_R)\tilde{Q}(\zeta,z^2/a^2) \nn\\
&= Z^{-1}_X(z^2\mu^2_R,\epsilon)\tilde{Q}(\zeta,z^2,\epsilon)\,.
\end{align}
As a result, we can relate $\tilde{Q}^X(\zeta,z^2\mu^2_R)$ to the $\overline{\rm MS}$ scheme by the conversion
\begin{align}
\tilde{Q}^X(\zeta,z^2\mu^2_R) =& \frac{Z_{\overline{\rm MS}}(\epsilon,\mu)}{Z_X(z^2\mu^2_R,\epsilon)}\, \tilde{Q}^{\overline{\rm MS}}(\zeta,z^2\mu^2)\nn\\
=& Z'_X(z^2\mu_R^2,\mu_R^2/\mu^2)\, \tilde{Q}^{\overline{\rm MS}}(\zeta,z^2\mu^2)\,,
\end{align}
where the regulator $\epsilon$ dependence is completely canceled out between $Z_{\overline{\rm MS}}$ and $Z_X$. The ratio $Z_X'$ can be calculated perturbatively in QCD, which was done in~\cite{Constantinou:2017sej} for several lattice schemes and the RI/MOM scheme. Thus the factorization formula we have proven in Sec.~\ref{sec:ope} still applies to $\tilde{Q}^X$ with a slight modification to the coefficient function,
\begin{align}
\tilde{Q}^X(\zeta,z^2\mu^2_R)
 = \int_{-1}^1 \!\! d\alpha \
   \mathcal{C}^X(\alpha,\mu_R^2/\mu^2,\mu^2 z^2)\, Q(\alpha\zeta,\mu)\,,
\end{align}
where the matching coefficient for the scheme $X$ is related to that of $\overline{\rm MS}$ by
\begin{align}
\mathcal{C}^X(\alpha,\mu_R^2/\mu^2,\mu^2 z^2) = Z'_X(z^2\mu_R^2,\mu_R^2/\mu^2)\, \mathcal{C}(\alpha,\mu^2 z^2) \,.
\end{align}
For the pseudo-PDF the modified result also involves this same coefficient
\begin{align}
\mathcal{P}^X(x,z^2\mu^2_R)
=&\int_{|x|}^1 {dy\over |y|}\
 \mathcal{C}^X\Bigl({x\over y},{\mu_R^2\over \mu^2},\mu^2 z^2\Bigr) q(y,\mu)
  \\
&+\int^{-|x|}_{-1} {dy\over |y|}\
 \mathcal{C}^X\Bigl({x\over y},{\mu_R^2\over \mu^2},\mu^2 z^2\Bigr) q(y,\mu)
 \,. \nn
\end{align}
Meanwhile, for the quasi-PDF we have,
\begin{align}\label{eq:rimomfac}
\tilde{q}_X\left(x,{\mu_R^2\over P_z^2}\right)
&\equiv \int \frac{d \zeta}{2 \pi}\: e^{ix \zeta}\:  \tilde{Q}^X \biggl(\zeta, \frac{{\mu}^2\zeta^2}{P_z^2}\biggr)
\\
&=\int_{-1}^1 {dy\over |y|}
\: C^X\Bigl(\frac{x}{y},\frac{\mu_R}{\mu},\frac{\mu}{|y| P^z}\Bigr)\, q(y,\mu)\,.\nn
\end{align}
Here the modified coefficient for the $X$ scheme is related to coefficient in the $\overline{\rm MS}$ scheme by
%
\begin{align}
 & C^X\Bigl(\frac{x}{y},  \frac{\mu_R}{\mu},\frac{\mu}{|y| P^z}\Bigr)
 \\
&\ =\!\! \int\!\! d\eta\ \bar{Z}'_X\Bigl(\eta^2, {\mu_R^2\over \mu^2}\Bigr)\:
 C\Bigl( \frac{x}{y} - \frac{\eta}{|y|}\frac{\mu_R}{P^z},{\mu\over |y|P^z}\Bigr)
 , \nn
\end{align}
where here $\bar{Z}'_X$ is defined by the Fourier transform
\begin{align}
\bar{Z}'_X\left( \eta^2,{\mu_R^2\over \mu^2}\right)
\equiv  \int {d\tau\over 2\pi}\ e^{i\eta \tau }\, Z'_X\Bigl(\tau^2,\frac{\mu_R^2}{\mu^2}\Bigr)
  \,.
\end{align}
Depending on the scheme $X$ we note that slightly modified definitions of $\bar{Z}'_X$ may be more appropriate.

One undesirable feature of  the $\overline{\rm MS}$ scheme for the renormalized \spcorr is that it does not have a smooth $z\to 0$ limit, and hence no simple connection with the fact that the local operator for $z=0$ is a conserved current. To avoid this one can simply make use of a different scheme that has a simple relation to $\overline{\rm MS}$, such as by adding $C_0(\mu^2z^2)$ to the $\overline{\rm MS}$ renormalization constant. This removes the offending $\ln(\mu^2 z^2)$ terms and yields a scheme with a smooth connection to the conserved current.

\begin{figure*}
	\centering
	\includegraphics[width=.49\textwidth]{images/Compare_qPDF_MSbar_Pzycut_TeX.pdf}
	\vspace{-0.3cm}
	\caption{
		The $\overline{\rm MS}$ scheme PDF $xf_{u-d}$ and the $\overline{\rm MS}$ quasi-PDF obtained from $x\, C^{\overline{\rm MS}}(p^z)\otimes f_{u-d}$, comparing results obtained with $p^z=yP^z$ and $p^z=P^z$.
	}
	\label{fig:quasi}
\end{figure*}

\begin{figure*}
	\centering
	\includegraphics[width=.49\textwidth]{images/Compare_pPDFvsPDF_TeX.pdf}
	\hspace{0.1cm}
	\includegraphics[width=0.49\textwidth]{images/Compare_pPDFvsPDF_z_TeX.pdf}	\\
	\vspace{-0.3cm}
	\caption{(Left) Comparison between the  PDF $xf_{u-d}$ and the pseudo-PDF  $x({\cal C}^{\overline{\rm MS}}\otimes f_{u-d})$ in the $\overline{\rm MS}$ scheme. The orange and blue bands indicate the results from varying the factorization scale $\mu=4\,{\rm GeV}$ by a factor of two. (Right) Same but now showing only central pseudo-PDF curves for different values of $z$. }
	\label{fig:pseudo}
\end{figure*}

Besides the $\overline{\rm MS}$ scheme, the quasi-PDF has also been defined with a transverse momentum cut-off~\cite{Xiong:2013bka,Lin:2014zya,Ma:2014jla,Alexandrou:2015rja} and in the RI/MOM scheme~\cite{Constantinou:2017sej,Alexandrou:2017huk,Chen:2017mzz,Green:2017xeu,Stewart:2017tvs}.  The RI/MOM scheme has attracted strong interest recently as it can be implemented nonperturbatively on the lattice, so we consider it as an explicit example of the above relations. In this scheme, the renormalization constant $Z_{\rm OM}$ is determined by imposing a condition on the \spcorr in an off-shell quark state,
\begin{align}
&\left.Z^{-1}_{\rm OM}\,
\tilde{Q}(\zeta=zp^z,z^2/a^2,-p^2a^2)\right|_{p^2=-\mu_R^2,p^z=p_R^z}
\nn\\
& \quad = \tilde{Q}^{(0)}_q(zp_R^z,z^2/a^2,z^2\mu_R^2) = e^{-izp_R^z}\,,
\end{align}
where $q$ denotes the quark state, $p^\mu$ is the external momentum, and ``$(0)$" in the superscript stands for the tree-level matrix element. As a result,
\begin{align}
Z_{\rm OM} &= Z_{\rm OM} (zp_R^z, z^2/a^2,a^2\mu_R^2)\,,\nn\\
Z'_{\rm OM}&=Z'_{\rm OM}(zp_R^z, z^2\mu_R^2, \mu_R^2/\mu^2)\,,
\end{align}
and here we define
\begin{align}
	&\bar{Z}'_{\rm OM}\left(\eta,{\mu_R^2\over (p_R^z)^2},{\mu_R^2\over \mu^2}\right)\nn\\
	&\equiv
   p_R^z \int {dz\over 2\pi}\ e^{i\eta p_R^z z}\ Z'_{\rm OM}(zp_R^z, z^2\mu_R^2, \mu_R^2/\mu^2)\,.
\end{align}
Then the matching coefficient in Eq.~(\ref{eq:rimomfac}) becomes
\begin{align}
& C^{\rm OM}\left(\frac{x}{y}, {\mu_R\over p_R^z},{\mu_R\over \mu}, {\mu\over yP^z}\right)
 \\
=& \int d\eta\ \bar{Z}'_{\rm OM}\left(\eta,{\mu_R^2\over (p_R^z)^2},{\mu_R^2\over\mu^2}\right) C\left({x\over y}-{\eta\over y}{p_R^z\over P^z}, {\mu\over |y|P^z}\right)
 \,. \nn
\end{align}
The choices of $\mu_R$ and $p_R^z$ are independent of $\mu$ and $P^z$, and $p_R^z=P^z$ was used in Refs.~\cite{Chen:2017mzz,Stewart:2017tvs}.

It should be noted that on the lattice, due to the breaking of chiral symmetry, the vector-like quark Wilson line operator $\tilde{O}_{\gamma^\mu}(z)$ can mix with the scalar operator $\tilde{O}_{\bf 1}(z)$, as has been discussed in Refs.~\cite{Constantinou:2017sej,Alexandrou:2017huk,Chen:2017mzz,Green:2017xeu,Chen:2017mie}. After considering the mixing effects, the same factorization formula can still be applied to the RI/MOM quasi-PDF from lattice QCD.
